\documentclass[../../main.tex]{subfiles}
\graphicspath{{../../image/}{../../figure/}}
\begin{document}

\section{经由原点处具锥奇异的熵的 $L^1$ 稳定性}\label{section:ch16-原点锥奇异熵l1稳定}

只有当初值属于 $\mathscr{C}_\omega$ 且具有生成满足\eqref{align:ch16-s03-e01}之解的性质时，才能应用\cref{theorem:ch16-1631}。本节与下一节的目标是证明：当源耗散时，类 $\mathscr{C}_\omega$ 相当丰富，能包含平衡态附近令人感兴趣的初值。

我们假设系统\eqref{align:ch16-s02-e01}配有一个熵--熵流对 $(\eta,q)$，其中 $\eta(U)$ 凸，从而可容解 $U$ 必须满足熵不等式
\begin{align}
\partial_t\eta(U(x,t))+\partial_xq(U(x,t))+\mathrm{D}\eta(U(x,t))P(U(x,t))\leqslant0.
\label{align:ch16-s05-e01}
\end{align}
如同第 5.5 节那样，称源 $P$ 相对于 $\eta$ 耗散，若熵产生非负：
\begin{align}
\mathrm{D}\eta(U)P(U)\geqslant0,\qquad U\in\mathscr{O}.
\label{align:ch16-s05-e02}
\end{align}

达到我们目标的最直接方式，是寻求一个凸熵 $\eta$，它展现出锥奇异 $\beta^{-1/2}|U|\leqslant\eta(U)\leqslant\beta^{1/2}|U|$（$U$ 在原点附近），并使源耗散\eqref{align:ch16-s05-e02}。事实上，在那种情形下，由\eqref{align:ch16-s05-e01}，$\mathscr{C}_\omega$ 必包含所有满足 $\|U_0\|_{L^1}\leqslant\omega/\beta$ 的初值 $U_0$。我们现在展示如何为系统原书(5.5.52)构造一个具有上述规格的熵函数，该系统在文献中一直充当松弛现象系统的范例，受制于次特征条件原书(5.5.54)。
此处我们假设原书(5.5.54)以严格不等式形式成立。

我们假设 $f(0)=0$，从而原点是平衡态。特征速度为 $\lambda=-a(u)$ 与 $\mu=a(u)$，相应的特征向量为
\begin{align}
R=\frac{1}{2a}\begin{pmatrix}1\\-a\end{pmatrix},\qquad S=\frac{1}{2a}\begin{pmatrix}1\\a\end{pmatrix}.
\label{align:ch16-s05-e03}
\end{align}
于是，由\eqref{align:ch16-s02-e03}，
\begin{align}
A=\frac{1}{2a}\begin{pmatrix}a+f'&-a+f'\\-a-f'&a-f'\end{pmatrix},
\label{align:ch16-s05-e04}
\end{align}
在 $u=0$ 处取值。由于次特征条件，$A$ 的主对角线元素为正，从而\eqref{align:ch16-s03-e02}满足。然而，由于 $\det A=0$，严格对角占优条件\eqref{align:ch16-s02-e04}对任何特征向量选择都不成立，因而\cref{theorem:ch16-1621}不适用。着眼于用\cref{theorem:ch16-1631}代替\cref{theorem:ch16-1621}，我们寻求一个定义在原点某邻域上的凸熵 $\eta(u,v)$，展现出锥奇异
\begin{align}
\beta^{-1/2}(|u|+|v|)\leqslant\eta(u,v)\leqslant\beta^{1/2}(|u|+|v|),
\label{align:ch16-s05-e05}
\end{align}
且产生非负：
\begin{align}
\eta_v(u,v)[v-f(u)]\geqslant0.
\label{align:ch16-s05-e06}
\end{align}

遵循第 12.2 节中的讨论，我们将把熵确定为 Riemann 不变量
\begin{align}
z=\int_0^u a(s)\,\mathrm{d}s-v,\qquad w=\int_0^u a(s)\,\mathrm{d}s+v
\label{align:ch16-s05-e07}
\end{align}
的函数 $\eta(z,w)$。

此处原书(12.2.2)约化为
\begin{align}
\eta_{zw}(z,w)+\kappa(z,w)[\eta_z(z,w)+\eta_w(z,w)]=0,
\label{align:ch16-s05-e08}
\end{align}
其中 $\kappa$ 表示 $(2a)^{-2}a'$ 表成 $(z,w)$ 的函数。此外，陈述 $\eta(u,v)$ 凸的原书(12.2.3)，此处取形式
\begin{align}
\eta_{zz}-\eta_{zw}\geqslant0,\qquad\eta_{ww}-\eta_{zw}\geqslant0.
\label{align:ch16-s05-e09}
\end{align}

局部平衡曲线 $v=f(u)$ 变成 $w=g(z)$，斜率为
\begin{align}
g_z=\frac{a+f'}{a-f'},
\label{align:ch16-s05-e10}
\end{align}
由次特征条件它为正。具有所需性质的熵 $\eta(z,w) $ 将在矩形
\begin{align}
\mathscr{R}=\{(z,w):-r<z<r,\ g(-r)<w<g(r)\},
\label{align:ch16-s05-e11}
\end{align}
上构造，其中 $r>0$ 待在下面指定。

我们固定一个 $(-\infty,\infty)$ 上的光滑函数 $\phi(z)$，满足 $\phi(0)=0$、$\phi_z(0)=0$，且对所有 $z\in(-\infty,\infty)$ 有 $|\phi_z(z)|\leqslant3$、$\phi_{zz}(z)\geqslant0$。在下文中，阶的符号 $O$ 将被理解为对一切具有上述规格的 $\phi$ 一致成立，特别地不依赖于 $\phi_{zz}$ 的大小。于是我们在 $\mathscr{R}$ 上把 $\eta(z,w)$ 构造为\eqref{align:ch16-s05-e08}在类空平衡曲线 $w=g(z)$ 上带 Cauchy 数据
\begin{align}
\eta(z,g(z))=\phi(z),\qquad-r<z<r,
\label{align:ch16-s05-e12}\\
\eta_w(z,g(z))-\eta_z(z,g(z))=0,\qquad-r<z<r
\label{align:ch16-s05-e13}
\end{align}
的解。

对\eqref{align:ch16-s05-e08}与原书(12.2.4)作了冗长但初等的估计后（见\cref{section:ch16-注记}所引文献），我们推得
\begin{align}
\eta(z,w)=\frac{1}{1+g_z(z)}\phi(z)+\frac{g_z(g^{-1}(w))}{1+g_z(g^{-1}(w))}\phi(g^{-1}(w))+O(z^2+w^2),
\label{align:ch16-s05-e14}\\
\eta_z(z,g(z))=\eta_w(z,g(z))=\frac{\phi_z(z)}{1+g_z(z)}=O(1),
\label{align:ch16-s05-e15}\\
\chi(z,w)\eta_{zz}(z,w)=\frac{\phi_{zz}(z)}{1+g_z(z)}+O(1),
\label{align:ch16-s05-e16}\\
\chi^*(z,w)\eta_{ww}(z,w)=\frac{\phi_{zz}(g^{-1}(w))}{g_z(g^{-1}(w))[1+g_z(g^{-1}(w))]}+O(1),
\label{align:ch16-s05-e17}
\end{align}
其中 $\chi(z,w)$ 与 $\chi^*(z,w)$ 是积分因子：
\begin{align}
\chi(z,w)=\exp\int_{g(z)}^w\kappa(z,\xi)\,\mathrm{d}\xi,\qquad\chi^*(z,w)=\exp\int_{g^{-1}(w)}^z\kappa(\zeta,w)\,\mathrm{d}\zeta.
\label{align:ch16-s05-e18}
\end{align}

我们取 $\phi(z)=\psi(z)+\gamma z^2$，其中 $\psi(0)=0$、$\psi_z(0)=0$，且对所有 $z\in(-\infty,\infty)$ 有 $|\psi_z(z)|\leqslant1$、$\psi_{zz}(z)\geqslant0$。正常数 $\gamma$ 取充分大，常数 $r$ 成比例地小，以使 $\gamma r<1$，从而对 $(z,w)\in\mathscr{R}$ 保证 $|\phi_z(z)|\leqslant3$ 与\eqref{align:ch16-s05-e07}都成立。于是 $\eta(u,v)$ 将是凸的。此外，由链式法则，$\eta_v=\eta_w-\eta_z$、$\eta_{vv}=\eta+\eta_{ww}-2\eta_{zw}$，从而\eqref{align:ch16-s05-e13}与\eqref{align:ch16-s05-e09}一起蕴涵\eqref{align:ch16-s05-e06}。

此时，保持 $r$ 与 $\gamma$ 如上固定，我们考虑一个 $\psi(z)$ 序列，它收敛到 Lipschitz 函数 $|z|$。相应熵序列然后收敛到 $\mathscr{R}$ 上某个 Lipschitz 函数 $\eta$，它是原书(5.5.52)的一个熵，作为 $(u,v)$ 的函数是凸的，并满足\eqref{align:ch16-s05-e06}。
此外，由\eqref{align:ch16-s05-e14}，
\begin{align}
\eta(z,w)=\frac{1}{1+g_z(z)}|z|+\frac{g_z(g^{-1}(w))}{1+g_z(g^{-1}(w))}|g^{-1}(w)|+O(z^2+w^2)+\gamma O(z^2+w^2).
\label{align:ch16-s05-e19}
\end{align}
作为最后一步，保持 $\gamma$ 如上固定，必要时进一步缩小 $r$ 的大小，以使\eqref{align:ch16-s05-e19}中线性增长的项控制二次阶的项，此时 $\eta$ 也满足\eqref{align:ch16-s05-e05}。

由于构造的 $\eta$ 满足\eqref{align:ch16-s05-e05}与\eqref{align:ch16-s05-e06}，熵不等式\eqref{align:ch16-s05-e01}蕴涵原书(5.5.52)的 $\mathscr{C}_\omega$ 包含所有满足
\begin{align}
\int_{-\infty}^{\infty}\{|u_0(x)|+|v_0(x)|\}\,\mathrm{d}x=\sigma,
\label{align:ch16-s05-e20}
\end{align}
其中 $\sigma\leqslant\beta^{-1}\omega$ 的初值。因此，当初值既满足\eqref{align:ch16-s05-e20}又满足
\begin{align}
TV_{(-\infty,\infty)}u_0(\cdot)+TV_{(-\infty,\infty)}v_0(\cdot)=\delta,
\label{align:ch16-s05-e21}
\end{align}
且 $\sigma$ 与 $\delta$ 充分小时，\cref{theorem:ch16-1631}在大范围建立了 $BV$ 解的存在性。

构造一个具有上述规格——特别地在原点处具锥奇异——的熵，似乎只对配有丰富熵族的系统可行，最显著的是只含两个平衡律的系统。对于规模更大、熵集合有限的系统，必须使用一种替代方法，它将在下一节发展。

\end{document}
